\subsection*{i.}
\addcontentsline{toc}{subsection}{Inciso i}
\textbf{Una institución financiera necesita almacenar millones de registros relacionados con clientes, cuentas y transacciones. Aunque la capacidad de almacenamiento ya no representa una limitación importante, la organización sigue invirtiendo recursos en el diseño adecuado de sus bases de datos. Explica por qué la simple acumulación de datos no garantiza información útil y analiza la importancia de aspectos como organización, estructura, integridad y seguridad para el éxito de los sistemas de información.}

Diseñar una base de datos útil para un propósito especifico, requiere que los datos almacenados sean útiles para resolver dicho propósito. De este modo, los datos almacenados en una base de datos deben, en primer lugar, estar asociados con la cuestión a resolver y en segundo lugar, particularmente en bases de datos a gran escala y en crecimiento como en el caso de una institución financiera, se debe tener la capacidad de manipular dichos datos con facilidad, para extraer la información necesaria y cumplir con su propósito.

Para obtener información importante de una base de datos, se debe tener la capacidad de establecer relaciones entre los datos. Una colección de datos sin relaciones entre si, es una colección dispersa y difícil de interpretar, particularmente para este caso, si se tiene una colección con los datos de clientes, cuentas y transacciones, pero no se tiene la capacidad de relacionarlos entre si, se pierde la capacidad de obtener mayor información del conjunto de datos almacenados, como las cuentas y las transacciones que puede realizar un cliente.

Realizar un diseño adecuado para el almacenamiento de los datos, permite extraer información fiable, además de otorgar seguridad sobre la informacion obtenida. Para el diseño se debe considerar:
\begin{itemize}
    \item \textbf{Organización eficiente: } Es necesario tener una forma de almacenamiento en la que se puedan manipular y revisar los datos de manera eficiente. Particularmente en una base de datos enorme y creciente como la de una institución financiera, con múltiples consultas y modificaciones, es indispensable tener la capacidad de operar con la base de datos de forma rápida y evitar retrasos con la información proporcionada a los clientes.
    \item \textbf{Estructura adecuada: } Como se mencionó anteriormente, una colección de datos por si mismos no tienen la capacidad de proporcionar información adecuada. Por ello, es necesario darles una estructura correcta, con la que sea fácil consultarlos, y relacionarlos correctamente entre si para extraer la información necesaria y cumplan con su propósito.
    \item \textbf{Integridad: } Esta cualidad, permite establecer reglas para que los datos almacenados sean coherentes y garantizar que la información obtenida a partir de las relaciones entre datos sea consistente. Para el caso de una institución financiera, es de suma importancia que los datos almacenados tengan sentido, pues la existencia de datos corruptos sobre cuentas y transacciones realizadas por los clientes, puede ocasionar pérdidas económicas en el futuro. Por ejemplo, al realizar una transacción, es indispensable que el cliente tenga en su cuenta el monto suficiente para realizarla, sin esta restricción, el cliente puede realizar transacciones sin restricciones y ocasionarle pérdidas a la institución.
    \item \textbf{Seguridad: } Es indispensable tener un control sobre las operaciones que se pueden realizar sobre la base de datos y evitar vulnerabilidades que conlleven una perdida de información. En particular, para una institución financiera, se vuelve indispensable, pues perder o filtrar información sobre sus clientes implicaría una perdida económica y de fiabilidad para la institución.
\end{itemize}
Una colección de datos almacenada sin un diseño adecuado es una base de datos vulnerable y propensa a errores futuros. Por esta razón, invertir en el diseño de la base de datos, es una inversión para evitar pérdidas en el futuro y para obtener información consistente, lo cual es fundamental en cualquier institución.
